\section{Introduction}
\label{sec:introduction}

Day-ahead electricity price forecasts support short-term electricity market
decisions \cite{weron2014electricity,lago2021forecasting}. In the coordinated
European day-ahead market, orders must be submitted before gate closure at
12:00 CET/CEST\footnote{All clock times in this thesis refer to Central European
Time or Central European Summer Time, the local time of the \mbox{DE-LU} market.}
on the day before delivery
\cite{epexSpotPowerMarket,epexSpotMarketCoupling}. A price forecast for this
setting must therefore be created before this deadline. The same timing rule
also matters in historical backtests, because using information that was not
yet observable at the forecast creation time would introduce data leakage.

Load, solar, and wind generation are central inputs for such price
forecasts. Demand affects scarcity, while solar and wind generation reduce the
residual demand that has to be served by conventional plants
\cite{pape2016fundamentals,do2021residualDemand}. Through the merit-order
effect, generation with low short-run marginal costs can displace more
expensive units and reduce spot prices \cite{sensfuss2008meritOrder}. For this
reason, empirical price forecasting models often include load, renewable
generation, weather, and calendar information
\cite{weron2014electricity,lago2021forecasting,maciejowska2021enhancing}.
Some studies bring weather information into the price model directly. For
example, \textcite{ludwig2015bigdata} use a broad set of German weather
variables, and \textcite{sgarlato2023weatherPredictions} use meteorological
forecasts such as wind, irradiation, cloud cover, and temperature to forecast
German electricity prices.

However, several price-forecasting studies for the Germany-Luxembourg bidding
zone (\mbox{DE-LU}) use day-ahead load forecasts published through the European Network of Transmission System Operators for
Electricity (\mbox{ENTSO-E}) Transparency Platform as price-model inputs
\cite{trebbien2023probabilistic,marcjasz2023distributional,uniejewski2025smoothing,hilger2024multivariate,eiser2026operational},
and several also use the \mbox{ENTSO-E}-published solar and wind forecasts
instead of direct weather variables
\cite{trebbien2023probabilistic,marcjasz2023distributional,hilger2024multivariate}.
These published forecasts compress weather-sensitive system conditions into
expected load, solar, and wind generation, so a price model that
relies on them implicitly assumes that they are accurate, observable in time,
and sufficiently complete. They are, however, effectively black boxes:
mandatory transparency data reported by system operators through undisclosed
models and information sets rather than reproducible specifications
\cite{euRegulation543_2013,entsoeTransparency}. How good these forecasts are compared
with forecasts that are transparent and reproducible from open data, and what
they are worth for price forecasting, can therefore only be judged by
benchmarking them.

Timing is the second problem, and it decides whether a published forecast can be
used at all. For load, the
\mbox{ENTSO-E} day-ahead forecast is public before gate closure and can enter
forecast runs conducted after its publication \cite{euRegulation543_2013}. The
\mbox{ENTSO-E} solar and wind forecasts, by contrast, are used as price-model
inputs in several \mbox{DE-LU} studies
\cite{trebbien2023probabilistic,marcjasz2023distributional,hilger2024multivariate}
but are not publicly observable before the 12:00 gate closure, with a
publication requirement no later than 18:00 on day \(d-1\)
\cite{kleinebrahm2026energyarena,entsoeTransparency,euRegulation543_2013}. For
an operationally valid backtest, each input must be observable at the forecast
creation time. These forecasts can therefore serve as quality benchmarks but
not as pre-gate-closure inputs.
This is the operational gap addressed here, which this thesis fills with solar
and wind forecasts generated from open data before gate closure.

Because admissibility depends on public observability at the forecast creation
time, the information set is itself a central design choice. This thesis defines a
strict operational information set for each hourly forecast creation time
from 07:00 to 12:00 on day \(d-1\). These sets become richer towards gate
closure as newer weather runs, more recent
realized system data, reserve-market auction results, and the public load
forecast become observable.

\phantomsection
\label{sec:research_questions}
These considerations lead to three research questions (RQ1 to RQ3):

\begin{description}
  \setlength{\itemsep}{0.15em}
  \setlength{\parsep}{0pt}
  \setlength{\topsep}{0.25em}
  \item[\textbf{RQ1:}]\hypertarget{rq:component_forecast_quality}{}
  How do fully transparent, open-data-based forecasts of load, solar, and wind generation compare with the corresponding
  \mbox{ENTSO-E} forecasts?

  \item[\textbf{RQ2:}]\hypertarget{rq:price_impact_generated_forecasts}{}
  How does an explicit representation of load and renewable-generation
  forecasts affect electricity price forecasts?

  \item[\textbf{RQ3:}]\hypertarget{rq:cutoff_sensitivity}{}
  How does electricity price forecast performance change with richer
  information sets over time?
\end{description}

The following hypotheses make these questions testable:

\begin{description}
  \setlength{\itemsep}{0.15em}
  \setlength{\parsep}{0pt}
  \setlength{\topsep}{0.25em}
  \item[\textbf{H1:}]\hypertarget{hyp:component_forecast_quality}{}
  The self-generated load forecast has lower or comparable forecast errors than the
  \mbox{ENTSO-E} load forecast, while the self-generated solar and wind forecasts
  have higher errors than the corresponding \mbox{ENTSO-E} forecasts.

  \item[\textbf{H2:}]\hypertarget{hyp:price_impact_generated_forecasts}{}
  Adding explicit self-generated load and renewable-generation forecasts improves
  electricity price forecast quality, relative to a baseline price model that
  uses the published \mbox{ENTSO-E} load forecast and represents renewable
  generation only through raw meteorological inputs.

  \item[\textbf{H3:}]\hypertarget{hyp:cutoff_sensitivity}{}
  Price forecast accuracy improves as the forecast creation time moves closer to
  gate closure and the admissible information set expands during the morning
  before delivery.
\end{description}

The intuition behind H1 is an asymmetry in information. For load, the public
forecast is available before gate closure \cite{euRegulation543_2013}, so
the self-generated forecast can correct it with newer weather runs and more
recent realized load. For solar and wind, no public forecast is available before
gate closure that could be corrected in this way. The self-generated forecasts must rely on open weather and capacity
data alone, whereas the \mbox{ENTSO-E} forecasts are produced by the system
operators with undisclosed models and inputs and may be issued as late as 18:00
on day \(d-1\) \cite{euRegulation543_2013,entsoeTransparency}, so higher errors
are expected for these two targets.

The central contribution of this thesis is a set of fully operational forecasting
models for load, solar, and wind generation based exclusively on open data,
referred to as operational boosted-tree forecasting (OBTF). The
models are released as part of an open pipeline that is ready for operational
use and that others can rerun and reuse. Their forecasts, together with the resulting price forecasts, are
submitted every day to the open Energy-Arena benchmark
\cite{kleinebrahm2026energyarena}, so their accuracy can be verified
independently and continuously, also beyond the test period. In contrast to
the \mbox{ENTSO-E} forecasts, whose methods are not disclosed, the inputs and
construction of these models are fully open. By construction, each model
is admissible before gate closure, so it can be used where the corresponding
\mbox{ENTSO-E} forecast cannot, as for solar and wind. The benchmarking in RQ1
further shows the OBTF load model to be even more accurate than the \mbox{ENTSO-E}
load forecast. The price forecast is one application that demonstrates their
operational value.

The remainder of the thesis is structured as follows. Chapter~\ref{sec:literature}
reviews related work, Chapter~\ref{sec:methodology} describes the data,
models, and evaluation design, Chapter~\ref{sec:results} presents the results,
and Chapters~\ref{sec:discussion} and~\ref{sec:conclusion} discuss and
conclude the thesis.
